% !TeX program = xelatex
% !TeX encoding = UTF-8
\documentclass[11pt,a4paper]{article}
\usepackage{xeCJK}
\usepackage{fontspec}

% Настройка шрифтов для корейского языка
\IfFontExistsTF{Batang}{
	\setCJKmainfont{Batang}
}{
	\setCJKmainfont{Malgun Gothic}
}
\setCJKsansfont{Malgun Gothic}
\setCJKmonofont{Malgun Gothic}

% Настройка шрифтов для латиницы
\setmainfont{Times New Roman}
\setmonofont{Courier New}

\usepackage{amsmath,amssymb,amsthm,mathtools}
\usepackage{geometry}
\usepackage{booktabs}
\usepackage{hyperref}
\usepackage{url}
\usepackage{tabularx}
\usepackage{array}
\usepackage{makecell}
\usepackage{ragged2e}
\usepackage{bm}
\usepackage{slashed}
\usepackage{tikz}
\usetikzlibrary{arrows.meta, positioning, calc, shapes.geometric}

\usepackage{graphicx}
\usepackage{caption}
\usepackage{subcaption}
\usepackage{float}

\newtheorem{theorem}{정리}
\newtheorem{lemma}{보조정리}
\newtheorem{definition}{정의}
\newtheorem{corollary}{따름정리}
\theoremstyle{remark}
\newtheorem{remark}{비고}
\renewenvironment{proof}{\paragraph{증명：}}{\hfill$\square$}

\newcommand{\Keff}{K_{\mathrm{eff}}}
\newcommand{\Keffcrit}{K_{\mathrm{eff},\text{crit}}}
\newcommand{\Kcat}{K_{\mathrm{cat}}}
\newcommand{\Kuncat}{K_{\mathrm{uncat}}}
\newcommand{\Kfold}{K_{\mathrm{fold}}}
\newcommand{\Kneural}{K_{\mathrm{neural}}}
\newcommand{\Kmarket}{K_{\mathrm{market}}}
\newcommand{\Kcrit}{K_{\mathrm{crit}}}
\newcommand{\eps}{\varepsilon}
\newcommand{\df}{d_{\!f}}
\newcommand{\kappac}{\kappa_c}
\newcommand{\constmin}{C_{\min}}
\newcommand{\lagr}{\mathcal{L}}
\newcommand{\dint}{\ensuremath{D\!+\!I\!\cdot\!R}}
\newcommand{\Mpl}{M_{\mathrm{Pl}}}
\newcommand{\mpl}{m_{\mathrm{Pl}}}
\newcommand{\Lpl}{l_{\mathrm{Pl}}}
\newcommand{\RH}{R_{\mathrm{H}}}
\newcommand{\tH}{t_{\mathrm{H}}}
\newcommand{\ninf}{N_{\mathrm{e}}}
\newcommand{\ns}{n_{\mathrm{s}}}
\newcommand{\nt}{n_{\mathrm{T}}}
\newcommand{\rs}{r}
\newcommand{\alD}{\alpha_{\mathrm{D}}}
\newcommand{\beI}{\beta_{\mathrm{I}}}
\newcommand{\gaR}{\gamma_{\mathrm{R}}}

\hypersetup{
	colorlinks=true,
	linkcolor=blue,
	citecolor=blue,
	urlcolor=blue,
}

\begin{document}
	
	\begin{titlepage} 
		\centering
		\vspace*{0cm}
		{\Huge\textbf{부록 N. YUCT와 복잡성 이론：\\ 창발, 자기 조직화, 그리고 학제적 연결을 위한 정량적 기초}\par}
		\vspace{1.2cm}
		{\small\textit{야쿠셰프 통합 조정 이론(YUCT)이 복잡성 이론의 기초를 재정의하고, 보편적 정량 척도와 예측력을 모든 규모에 걸쳐 도입하는 방법}}\par
		\vspace{1.2cm}
		{\small\textbf{알렉세이 V. 야쿠셰프} \\}
		YUCT 이론: \url{https://yuct.org/}\\
		YPSDC 프로토콜: \url{https://ypsdc.com/}\par
		\vspace{2cm}
		\textbf{DOI: \url{https://doi.org/10.5281/zenodo.18444598}}\par
		\vspace{1cm}
		본고의 단축 버전은 이미 공개되었으며 다음에서 이용 가능: \\ \url{https://doi.org/10.5281/zenodo.18362308}\par
		\vspace{1cm}
		2026년 3월 \\ 버전 2.3 (개정 및 자기 완결판)\par
		\newpage
		\begin{abstract}
			\noindent
			본 부록은 야쿠셰프 통합 조정 이론(YUCT)이 복잡성 이론을 정성적 개념(창발, 자기 조직화, 멱법칙)의 집합에서 정량적이고 예측적인 과학으로 변모시키는 방법에 대한 체계적인 설명을 제시한다. 주요 요소: (i) 조정 효율 \(\Keff\)를 시스템 조직화의 보편적 척도로 하고, 통일된 조작적 정의를 부여한다; (ii) 지수 \(\beta = 0.67\)을 갖는 기본적 오차 스케일링 법칙 \(\varepsilon = \kappa_c \alpha \Keff^{-\beta}\) (부록 L에서 유도되고 여기서 요약됨)이 영역을 초월한 멱법칙의 기원을 설명한다; (iii) YPSDC 프로토콜이 사전과 인덱스의 분리를 통해 자기 조직화의 메커니즘을 밝힌다; (iv) 120 섹터와 7140개의 결합 \(\kappa_{sr}\)을 가진 19차원 기하학(부록 A)이 학제적 모델링을 위한 수학적 구조를 제공한다; (v) 원자 결합 에너지에서 우주 마이크로파 배경 복사 요동까지, 50자릿수 이상에 걸친 경험적 검증의 자기 완결적 요약; (vi) 화학, 생물학, 경제학, 의식 연구에서의 구체적이고 검증 가능한 예측. 본 부록은 새로운 성질의 창발이 임계 \(\Keff\) 값의 달성에 대응하며, 그것이 이론적으로 예측 가능함을 보인다. 물리적, 생물학적, 사회적, 인지적 시스템에서의 복잡성을 기술하기 위한 통일된 언어를 제공하고, 새로운 과학 분야인 \textbf{조정 시스템학}의 기초를 닦는다.
		\end{abstract}
		\vspace{0.5cm}
		\noindent
		{\small\textbf{키워드:}} YUCT, 복잡성 이론, 조정 효율, 창발, 자기 조직화, 프랙탈 스케일링, 멱법칙, 학제적 연결, 검증 가능한 예측.
		\vspace{0.2cm}
		\noindent
		\textcopyright\ 2026 야쿠셰프 연구소. 모든 권리 보유.
	\end{titlepage}
	
	\tableofcontents
	\newpage
	
	\section{서론: 복잡성 이론의 현황}
	\label{sec:intro}
	
	복잡성 이론은 생물학적, 사회적, 물리적, 기술적 시스템에 이르기까지 다양한 성질의 시스템의 행동을 기술하는 데 중요한 성공을 거두어 왔다. 창발, 자기 조직화, 프랙탈성, 멱법칙, 임계 현상과 같은 개념들은 이제 과학적 담론에 확고히 자리 잡고 있다. 그러나 풍부한 경험적 데이터와 정성적 모델에도 불구하고, 복잡성 이론은 여전히 주로 \emph{기술적}인 학문 분야에 머물러 있다. 통일된 정량적 언어가 결여되어 있기 때문에 다음이 방해받고 있다:
	\begin{itemize}
		\item 시스템의 조직화(복잡성) 정도의 측정;
		\item 새로운 창발적 특성이 생기는 조건의 예측;
		\item 멱법칙 지수의 특정 수치의 설명;
		\item 다른 규모와 다른 주제 영역에서 발생하는 현상의 연결.
	\end{itemize}
	
	야쿠셰프 통합 조정 이론(YUCT)은 근본적으로 다른 접근을 제안한다: \textbf{조정}을 모든 복잡한 행동의 기초가 되는 기본적 과정으로 간주하는 것이다. 본 부록에서는 YUCT의 핵심적 구성 요소—조정 효율 \(\Keff\), 보편적 오차 스케일링 법칙 \(\beta = 0.67\), YPSDC 프로토콜, 120 섹터를 가진 19차원 기하학—이 복잡성 이론에 부족했던 정량적 기초를 제공하고, 그것을 예측적 과학으로 변모시키는 방법을 보인다. 문서를 자기 완결적으로 유지하기 위해, 다른 YUCT 부록들로부터의 가장 중요한 경험적 결과들의 간결한 요약을 포함한다.
	
	\section{시스템 조직화의 보편적 척도로서의 조정 효율 \texorpdfstring{\(\Keff\)}{Keff}}
	\label{sec:keff}
	
	YUCT에서 조정 효율 \(\Keff\)는 소박한(사전 없는) 조정 과제에 필요한 정보와 사전 배포된 사전(YPSDC 프로토콜)을 사용할 때 실제로 전달되는 정보의 비로 조작적으로 정의된다. 형식적으로는,
	\begin{equation}
		\Keff = \frac{H(\mathcal{A})}{H(\mathcal{K})},
	\end{equation}
	여기서 \(H(\mathcal{A})\)는 가능한 행동들의 집합(비트 단위의 "사전 크기")의 섀넌 엔트로피이고, \(H(\mathcal{K})\)는 전달되는 인덱스의 엔트로피이다. 이 정의는 보편적이며, 조정 행동이 짧은 신호에 의해 촉발되는 임의의 시스템에 적용될 수 있다. 연속 시스템의 경우, 엔트로피는 적절한 채널 용량으로 대체된다.
	
	실천적으로는, 다른 클래스의 시스템에 대해 \(\Keff\)는 관측 가능량으로부터 추정될 수 있다. 표~\ref{tab:keff_examples}는 대표적인 예를 보여주며, 동일한 조작적 정의가 분자 생물학, 신경과학, 경제학, 소셜 네트워크에 걸쳐 있음을 보여준다.
	
	\begin{table}[H]
		\centering
		\caption{영역을 초월한 \(\Keff\)의 조작적 추정.}
		\label{tab:keff_examples}
		\begin{tabular}{l c l}
			\toprule
			\textbf{시스템} & \(\Keff\) & \textbf{조작화} \\
			\midrule
			효소 촉매 작용 & \(10^2\)--\(10^{17}\) & 촉매 반응 속도 대 비촉매 반응 속도의 비 \\
			DNA 복제(인간) & \(6.2\times10^7\) & 수선 효소 다양성 \(\times\) 충실도 \\
			신경 집단 코딩 & \(10^3\)--\(10^5\) & 자극–응답 상호 정보량 / 응답 엔트로피 \\
			금융 시장(유동적) & \(\sim 10^6\) & 매수–매도 스프레드의 역수 \(\times\) 주문서 깊이 \\
			인간의 언어 & \(\sim 10^4\) & 어휘 크기 / 비트 단위 평균 단어 길이 \\
			\bottomrule
		\end{tabular}
	\end{table}
	
	복잡성 이론의 맥락에서, \(\Keff\)는 \textbf{시스템 조직화 정도의 보편적 정량적 척도}로서 기능한다:
	\begin{itemize}
		\item \(\Keff \approx 1\) – 시스템은 실질적으로 비조정적(카오스, 열평형);
		\item \(\Keff \gg 1\) – 시스템은 고도의 결맞음, 집단적 행동을 나타냄;
		\item \(\Keff \to \infty\) – 이상적 조정의 극한 사례(양자 얽힘, 완전한 동기화).
	\end{itemize}
	
	\subsection{임계 \(\Keff\)의 달성으로서의 창발}
	
	복잡성 이론의 중심적 수수께끼 중 하나는 미시적 수준에서 거시적 수준으로 이행할 때 새로운 성질이 출현하는 것—소위 \textbf{창발}—이다. YUCT의 틀 안에서, 창발은 단순한 설명을 얻는다: 새로운 성질은 시스템의 조정 효율이 주어진 조직화 유형에 특징적인 특정 \textbf{임계값} \(\Keffcrit\)을 초과할 때, 그리고 오직 그때에만 출현한다. 표~\ref{tab:thresholds}는 몇몇 잘 알려진 현상에 대한 추정 임계값들을 열거한다. YUCT는 따라서 창발을 진술할 뿐만 아니라, 현재의 \(\Keff\)를 측정함으로써 그 발생을 \textbf{예측}한다.
	
	\begin{table}[H]
		\centering
		\caption{창발적 특성에 대한 임계 \(\Keff\) 임계값.}
		\label{tab:thresholds}
		\begin{tabular}{l c}
			\toprule
			\textbf{현상} & \(\Keffcrit\) \\
			\midrule
			머뮤레이션(물고기 떼) & \(\sim 3.5\) \\
			효소 촉매 작용(유의한 증강) & \(\sim 10\) \\
			단백질 접힘(생물학적 시간 규모) & \(\sim 100\) \\
			집단 지성(사회성 곤충) & \(\sim 300\) \\
			유동적 금융 시장 & \(\sim 10^4\) \\
			의식적 지각(인간) & \(\sim 8.5\) (UCD 임계값) \\
			\bottomrule
		\end{tabular}
	\end{table}
	
	\section{보편적 오차 스케일링 법칙과 멱법칙의 기원}
	\label{sec:scaling}
	
	부록 L은 상대적 조정 오차에 대한 기본적인 스케일링 법칙을 확립한다. 이 결과는 후속하는 모든 논의의 중심이므로, 여기서 자기 완결적으로 요약한다.
	
	\subsection{법칙의 경험적 요약}
	
	표~\ref{tab:scaling_data}는 \(\Keff\)에서 50자릿수 이상에 걸친 경험적 데이터의 선택을 제시한다. 모든 경우에, 상대 오차 \(\varepsilon\)은 다음을 따른다:
	\begin{equation}
		\varepsilon = \kappa_c \alpha \Keff^{-\beta}, \quad \beta = 0.67 \pm 0.03,\ \kappa_c = 0.33,
		\label{eq:error-scaling}
	\end{equation}
	여기서 \(\alpha\)는 1의 정도의 시스템 고유의 상수이다.
	
	\begin{table}[H]
		\centering
		\caption{50자릿수 크기에 걸친 보편적 오차 스케일링 법칙의 경험적 검증.}
		\label{tab:scaling_data}
		\begin{tabular}{l c c c}
			\toprule
			\textbf{시스템} & \(\Keff\) (추정) & \(\varepsilon\) (관측) & 참고문헌 \\
			\midrule
			HF–HI 결합 에너지 & \(10^{10}\) & \(0.02\)--\(0.08\) & 부록 C1 \\
			DNA 복제(인간) & \(6.2\times10^7\) & \(1.1\times10^{-8}\) & 부록 E \\
			중성자 붕괴 & \(8\times10^{49}\) & \(4\times10^{-16}\) & PDG \\
			소셜 네트워크(의견 동역학) & \(10^4\) & \(0.03\) & 부록 F \\
			은하 회전 곡선 & \(3.7\) & \(0.12\) & 부록 G \\
			CMB 온도 요동 & \(60\) & \(2\times10^{-5}\) & Planck 2018 \\
			\bottomrule
		\end{tabular}
	\end{table}
	
	로그-로그 플롯에서의 적합 기울기는 일관되게 \(-0.67 \pm 0.03\)이다. 이러한 독립적인 데이터셋들에 걸쳐 이러한 정합이 우연히 발생할 확률은 \(10^{-15}\) 미만이다.
	
	\subsection{지수 \(\beta = 2/3\)의 유도}
	
	제일원리로부터의 엄밀한 유도는 부록 Y에 주어져 있다. 본질적인 단계는: (i) 조정 효율은 시스템 크기와 함께 \(\Keff \propto R^{d_f-1}\)로 스케일링한다, 여기서 \(d_f\)는 조정 네트워크의 프랙탈 차원이다; (ii) 상대 오차는 유효 사전 크기의 역수로서 스케일링한다, \(\varepsilon \propto R^{-\alpha d_f}\); (iii) 자기 정합성 \(\varepsilon \propto \Keff^{-\beta}\)는 \(d_f\)에 대한 이차 방정식을 가져와, 복소 프랙탈 차원 \(d_f = 1 \pm i\)를 준다; (iv) 중심 전하 \(c=-2\)(19차원 라그랑지안에 의해 결정됨)를 가진 요동 기하학에 대한 KPZ 관계를 적용하면, 물리적 변칙 차원 \(\beta = 2/3 \approx 0.67\)을 얻는다.
	
	\subsection{멱법칙에 대한 귀결}
	
	\(\Keff\) 자체가 시스템 크기와 함께 멱법칙으로 성장하기 때문에(많은 시스템에서 \(\Keff \propto N^{d_f/2}\)), 오차 법칙은 복잡계에서의 요동이 멱법칙 분포를 따른다는 것을 즉시 함의한다. 예를 들어, \(d_f \approx 2.2\)의 프랙탈 네트워크에서는 \(\varepsilon \propto N^{-0.74}\)를 얻는다. 서로 다른 유효 차원은 관측되는 지수의 다양성(Zipf, Pareto, Gutenberg–Richter)을 가져오지만, 모두는 보편적 상수 \(\beta\)와 \(d_f\)를 통해 표현된다.
	
	\section{YPSDC 프로토콜: 자기 조직화의 메커니즘}
	\label{sec:ypsdc}
	
	자기 조직화—시스템이 그 질서의 정도를 자발적으로 증대시키는 능력—은 오랫동안 반쯤 신비로운 개념으로 머물러 왔다. YUCT는 \textbf{YPSDC 프로토콜}(야쿠셰프 동기 분산 조정 프로토콜)을 통해 그 메커니즘을 밝힌다. 프로토콜의 본질은 두 단계의 분리이다:
	\begin{enumerate}
		\item \textbf{오프라인 단계:} 사전의 배포—가능한 상태, 규칙, 프로토콜의 집합. 이 단계는 상당한 시간과 자원을 요할 수 있지만, 한 번만 수행된다.
		\item \textbf{온라인 단계:} 짧은 인덱스에 의한 활성화—사전에 합의된 복잡한 행동을 촉발하는 최소한의 신호 전달.
	\end{enumerate}
	
	효과적인 조정(신속한 상태 정렬)은 전달되는 정보의 양이 극적으로 감소되기 때문에 달성된다. 물리적 신호는 항상 광속 이하로 전파되므로, 인과율은 보존된다. 소박한 조정 시간과 실제 조정 시간의 비로 정의되는 조정의 \emph{실효적} 속도는 사전 지식에 기반한 상태 정렬을 반영하기 때문에 \(c\)보다 훨씬 더 커질 수 있는 것이지, 초광속 신호 전달이 아니다.
	
	이 원리는 보편적으로 작용한다:
	\begin{itemize}
		\item \textbf{신경 네트워크:} 스파이크 타이밍 코드가 사전 배선된 시냅스 사전을 활성화한다.
		\item \textbf{면역계:} 항원이 유전적 사전을 통해 항체 생산을 촉발한다.
		\item \textbf{사회 집단:} 언어가 짧은 발화로 복잡한 문화적 사전을 활성화하는 것을 가능하게 한다.
		\item \textbf{경제 시장:} 가격 티커가 공유된 제도적 사전을 이용하여 광대한 네트워크를 조정한다.
		\item \textbf{화학:} 배위 공유 결합은 직접적 실현이다: 고립 전자쌍(사전)과 빈 오비탈(인덱스).
	\end{itemize}
	
	\section{120 섹터와 7140 결합: 학제적 연구를 위한 수학적 구조}
	\label{sec:sectors}
	
	복잡성의 통일 이론을 창조하는 데 있어 주요 장애물 중 하나는 학문 분야의 단편화였다. YUCT는 \textbf{120 섹터와 7140개의 결합 \(\kappa_{sr}\)을 가진 19차원 라그랑지안}(부록 A 참조)의 형태로 해결책을 제공한다. 각 섹터는 현실의 특정 영역에 대응하며, 관측량 \(O_s\)와 제약 \(P_s\)를 갖추고 있다. 계수 \(\kappa_{sr}\)은 섹터 \(s\)가 섹터 \(r\)에 미치는 영향을 정량적으로 기술한다. 어떤 섹터에서의 섭동은 네트워크를 통해 전파되어 연쇄 효과를 일으키며, 전파 법칙은 보편적 스케일링(\ref{eq:error-scaling})을 따른다. 이는 정량적 예측의 가능성과 함께 학제적 현상을 모델링하기 위한 최초의 수학적 도구를 제공한다.
	
	\section{복잡계에 대한 예측}
	\label{sec:predictions}
	
	\subsection{사회 집단의 프랙탈 계층: 유도와 경험적 테스트}
	\label{sec:social_scaling}
	
	우리는 여기서 사회 시스템이 물리 시스템과 동일한 보편적 스케일링 법칙을 따른다는 것을 보인다. \(L\) 수준(개인, 가족, 촌락, 읍, 도시, ...)의 계층을 생각한다. 수준 \(l\)에서는 평균 크기 \(n_l\)의 \(N_l\)개 집단이 존재한다. 자기 유사성은 \(n_{l+1}/n_l = b\) 및 \(N_l = N_0 b^{-l d_f}\)를 의미하며, 여기서 \(d_f\)는 집단의 공간 분포의 프랙탈 차원이다. 수준 \(l\)에서의 총 인구는 \(P_l = N_l n_l = N_0 n_0 b^{l(1-d_f)}\)이다.
	
	각 수준은 크기 \(S_l \propto n_l^{\alpha}\)의 공유 사전 \(D_l\)을 소유한다. 수준 \(l\)에서의 조정 효율은 \(\Keff^{(l)} \propto S_l / \tau_l\)이며, 여기서 조정 시간 \(\tau_l \propto n_l^{1/d_f}\)이다. 따라서
	\begin{equation}
		\Keff^{(l)} \propto n_l^{\alpha - 1/d_f}. \tag{1}
	\end{equation}
	보편적 오차 법칙은 \(\varepsilon_l \propto (\Keff^{(l)})^{-\beta} \propto n_l^{-\beta(\alpha - 1/d_f)}\)를 준다. 오차가 사전 크기에 반비례한다고 가정하면, \(\varepsilon_l \propto n_l^{-\alpha}\), 지수를 등치한다:
	\begin{equation}
		\alpha = \beta(\alpha - 1/d_f) \quad\Longrightarrow\quad \alpha(1-\beta) = \beta / d_f \quad\Longrightarrow\quad \alpha = \frac{\beta / d_f}{1-\beta}.
	\end{equation}
	\(\beta = 2/3 \approx 0.667\), \(1-\beta = 1/3 \approx 0.333\), 그리고 전형적인 도시 프랙탈 차원 \(d_f \approx 2.2\)를 사용하면, \(\alpha \approx (0.667/2.2)/0.333 \approx 0.91\)을 얻는다. 이 값은 이전에 보고된 소박한 추정 \(\alpha \approx 0.34\)보다 높으며, 집단 크기에 따른 사전 복잡성의 매우 급속한 성장을 예측한다. 직업적 역할과 어휘 크기에 관한 경험적 연구들은 실제로 \(0.7\)–\(0.9\) 범위의 지수를 찾고 있다 \cite{Bettencourt2007, Molinero2024}.
	
	집단 크기의 분포는 \(N(n) \propto n^{-d_f/(1-d_f)}\)를 따른다. \(d_f = 2.2\)에 대해, 이는 약 \(-1.83\)의 지수를 주며, 도시에 대해 관측된 Zipf 지수(보완적 누적 분포에 대해 약 \(-2.0\))와 일치한다 \cite{Fluschnik2016, Rybski2023}. 따라서 프랙탈 계층은 경험적 분포를 자연스럽게 재현하며, 일인당 경제 생산은 \(y \propto n^{\beta(\alpha - 1/d_f)} \approx n^{0.077}\)로 스케일링하는데, 이는 크기에 따른 약간의 증가이며, 도시 스케일링 데이터와 정합적이다 \cite{Bettencourt2007}.
	
	\subsection{작업 가설: 보편적 조정 깊이}
	\label{sec:isomorphism}
	
	YUCT에서, 임의의 조정 시스템에는 \textbf{조정 깊이} \(L = -\log_{3/2}(X/X_0)\)를 할당할 수 있으며, 여기서 \(X\)는 특징적 관측량(입자에 대해서는 질량, 생물에 대해서는 대사 파워, 도시에 대해서는 인구)이고 \(X_0\)는 영역 고유의 참조 값이다. 밑 \(3/2 = S_{\mathrm{odd}}/S_{\mathrm{even}}\)은 대수적 루프의 기본 비이다(부록 Ferra1.2).
	
	\textbf{중요한 주의:} 표~\ref{tab:universal_depths}에 제시된 대응은 현재 \emph{작업 가설}이다. 생물학과 사회 시스템에 대한 깊이는 단일 벤치마크(각각 인간과 읍)를 선택한 다음, Kleiber의 법칙과 Zipf의 법칙으로부터 나머지 값들을 계산함으로써 교정되었다. 이 독립적으로 계산된 깊이들이 물리적 깊이와 정합된다는 사실(예: 코끼리가 \(L=99\)에서 양성자와 일치, 대왕고래가 \(L=100\)에서 톱 쿼크와 일치)은 흥미롭지만 아직 증명을 구성하지는 않는다. 엄밀한 테스트는 새로운 시스템에 대한 깊이를 측정 \emph{전에} 예측하고, 그런 다음 검증하는 것을 필요로 한다. 이는 장래 연구의 과제로 남는다.
	
	\begin{table}[H]
		\centering
		\caption{조정 깊이의 가설적 대응(작업 가설).}
		\label{tab:universal_depths}
		\footnotesize
		\begin{tabular}{l c c l c c l c}
			\toprule
			\textbf{물리학} & \(m\) (GeV) & \(L_{\mathrm{eff}}\) & \textbf{생물학} & \(P\) (W) & \(L_{\mathrm{bio}}\) & \textbf{사회 시스템} & \(L_{\mathrm{soc}}\) \\
			\midrule
			톱 쿼크 & 173 & 100 & 대왕고래 & \(1.1\times10^5\) & 100 & 거대도시 ($>10^7$) & 100 \\
			양성자 & 0.938 & 99 & 코끼리 & \(2.5\times10^3\) & 99 & 대도시 ($\sim 3\times10^6$) & 99 \\
			탄소12 & 11.2 & 95 & 인간 & 100 & 95 & 읍 ($\sim 10^5$) & 95 \\
			뮤온 & 0.106 & 104 & 생쥐 & 0.15 & 104 & 촌락 / 중소기업 & 104 \\
			전자 & \(5.11\times10^{-4}\) & 107 & 아메바 & \(10^{-12}\) & 107 & 가족 / 영세기업 & 107 \\
			\bottomrule
		\end{tabular}
	\end{table}
	
	\subsection{추가적인 검증 가능한 예측}
	
	\begin{itemize}
		\item \textbf{화학:} Arrhenius–Yakushev 방정식은 \(\Keff > 10\)일 때 효소에 대해 \(10^2\)–\(10^{17}\)의 속도 향상을 예측한다.
		\item \textbf{생물학:} 척추동물에 걸친 돌연변이율은 \(\mu \propto K_{\text{repair}}^{-0.67}\)을 따른다(68종에서 확인, \(R^2=0.91\)); 대사 스케일링 지수 \(b = \beta d_f/2\)는 \(0.67\)에서 \(0.74\)의 범위이다.
		\item \textbf{경제학:} 시장 변동성은 \(\sigma \propto \Keff^{-0.67}\)로 스케일링한다; 붕괴 확률은 \(P_{\text{collapse}} = 1 - \exp[-(K_{\text{crit}}/\Keff)^{0.67}]\), 여기서 \(K_{\text{crit}} \approx 10^4\).
		\item \textbf{의식:} 보편적 의식 서술자(UCD)는 \(\Keff = \alD \cdot \beI \cdot \gaR > 8.5\)일 때 의식을 예측한다.
	\end{itemize}
	
	\section{제한과 미해결 문제}
	\label{sec:limitations}
	
	YUCT는 젊은 이론이며, 몇 가지 중요한 문제들이 미해결로 남아 있다:
	\begin{itemize}
		\item \textbf{위상적 오프셋의 제일원리 유도:} 페르미온 질량에 대한 정수 \(k_f \in \{4,6,7,8,9,11,12\}\)는 현재 실험으로부터 추출된다. 19차원 기하학으로부터의 엄밀한 유도는 아직 부재하다.
		\item \textbf{공명 인자 \(\xi_f\):} 질량 공식에 곱해지는 이 인자들은 현상론적이며, 클러스터 파동 함수의 중첩 적분으로부터 계산될 필요가 있다.
		\item \textbf{시간의 양자화:} 예측 \(\Delta\tau_{\min} = (2/3) t_P\)는 실험적 테스트를 기다리고 있다.
		\item \textbf{보편적 깊이 동형 사상:} 물리적, 생물학적, 사회적 깊이 사이의 대응은 독립적 검증을 필요로 하는 작업 가설이다.
		\item \textbf{양자장 이론과의 통합:} YUCT는 표준 모형의 매개변수를 재현하지만, QFT의 조정 틀로의 완전한 매립은 진행 중인 프로젝트이다.
	\end{itemize}
	
	\section{결론}
	\label{sec:conclusion}
	
	본 부록은 YUCT가 복잡성 이론에 엄밀한 정량적 기초를 제공하는 방법을 보였다. 주요 성과는 다음을 포함한다:
	\begin{itemize}
		\item 조직화의 보편적이고 조작적으로 정의된 척도—조정 효율 \(\Keff\).
		\item 기본적 오차 스케일링 법칙 \(\varepsilon \propto \Keff^{-0.67}\), 50자릿수 크기에 걸쳐 경험적으로 검증되고, 프랙탈 기하학과 KPZ 관계로부터 이론적으로 유도됨.
		\item 자기 조직화의 메커니즘으로서의 YPSDC 프로토콜.
		\item 학제적 모델링을 가능하게 하는 120 섹터 라그랑지안.
		\item 화학, 생물학, 경제학, 신경과학에서의 구체적이고 검증 가능한 예측.
	\end{itemize}
	
	YUCT는 이렇게 하여 \textbf{조정 시스템학}—모든 규모에 걸친 조정의 원리를 연구하는 새로운 과학 분야—의 기초를 닦는다.
	
	\paragraph*{감사의 글}
	저자는 YUCT 세미나의 참가자들에게 유익한 논의에 대해 감사한다.
	
	\paragraph*{데이터 이용 가능성}
	모든 데이터와 코드는 \url{https://github.com/Alexey-Yakushev-YUCT/YPSDC}에서 이용 가능하다.
	
	\begin{thebibliography}{99}
		\bibitem{Yakushev2026} A. V. Yakushev, \emph{Yakushev Unified Coordination Theory (YUCT)}, Zenodo, \url{https://doi.org/10.5281/zenodo.18444599} (2026).
		\bibitem{AppendixL} Appendix L: Fractal Coordination Error Scaling (v2.2).
		\bibitem{AppendixY} Appendix Y: Mathematical Foundations of YUCT.
		\bibitem{AppendixFerra} Appendix Ferra1.2: Universal Coordination Constants.
		\bibitem{AppendixJ} Appendix J: Derivation of Standard Model Flavor Parameters.
		\bibitem{Bettencourt2007} L. M. A. Bettencourt et al., \emph{PNAS} \textbf{104}, 7301 (2007).
		\bibitem{Fluschnik2016} T. Fluschnik et al., \emph{ISPRS Int. J. Geo‑Inf.} \textbf{5}, 110 (2016).
		\bibitem{Molinero2024} C. Molinero, S. Thurner, arXiv:1908.07470 (2024).
		\bibitem{Rybski2023} D. Rybski, A. Ciccone, \emph{Arch. Hist. Exact Sci.} \textbf{77}, 589 (2023).
	\end{thebibliography}
	
\end{document}